\subsection{From contact slack to a conserved phase resource}
\label{ms:conserved}\label{ms:fourier:afcore:sec:conserved}

The physical defect measures how far the gap and third-moment geometry
are from the extremal two-point configuration. To bound that defect, we
integrate the contact slack in two ways. Integrating over support gaps
reveals $J$ and a nonnegative leakage $L$. Integrating against the
zero-bias law instead turns the same slack into a quadratic-hinge
expectation for the sum. A compensated Fourier formula then compares
that expectation with its Gaussian and cubic terms. Equating the two
expressions gives an exact identity in which score excess, defect,
phase displacement and leakage share one remainder.


We now state precisely the contact input needed for the identity.
Fix an integer $n\ge2$ and a common iid marginal $P\in\mathcal P$.
Write $S_j=X_1+\cdots+X_j$ and $z=H/\sqrt n$,
\[
 F=\Pp(S_n\le H),\quad
 C=\frac{\sqrt n}{\rho}\{F-\Phi(z)\}>0,\quad
 \beta=\E[X|X|],\quad \gamma=\E X^3,\quad \varphi=\phi(z).
\]
Suppose that the fixed-index contact argument has supplied
\begin{equation}\label{ms:fourier:afcore:eq:contact-input}
 h_P(y):=\Pp(S_{n-1}\le H-y)\le Q_H(y),
 \qquad h_P(y)=Q_H(y)\quad(y\in\operatorname{supp}P),
\end{equation}
where
\begin{align}
 nQ_H(y)={}&nF+\frac{H\varphi+C\rho}{2\sqrt n}
 -\left(\sqrt n\varphi+\frac{3C\beta}{\sqrt n}\right)y
 -\frac{H\varphi+3C\rho}{2\sqrt n}y^2+\frac{C}{\sqrt n}|y|^3.
 \label{ms:fourier:afcore:eq:contact-polynomial}
\end{align}
At an unrestricted maximizing positive closed-CDF witness, these
relations follow from Lemma~\ref{ms:contact}. They will be the precise
contact hypotheses throughout the derivation.

Because $C>0$, the piecewise cubic function grows to $+\infty$ on either support ray.
Equality in \eqref{ms:fourier:afcore:eq:contact-input}, with $0\le h_P\le1$,
therefore gives $|X|\le R$ for some finite $R$.
Only after this observation do we introduce
$M_4=\E X^4$ and $M_5=\E|X|^5$.

\subsubsection{Auxiliary gap integrals and leakage}

Use the zero-bias density $\tau$, gap functional $J$ and physical defect
$D$ of \eqref{ms:tau-moments} and
\eqref{ms:physical-defect-definition}. Lemma~\ref{ms:physical-defect}
gives $D\ge0$. To evaluate the contact slack we need two more weighted
gap integrals and one nonnegative leakage term. Let $\nu$ be the law of
$H-S_{n-1}$ and set
\begin{align}
 J_1&=\sum_{(a,b)}\tau_{a,b}\int_a^b(t-a)t\,dt,\qquad
 J_2=\sum_{(a,b)}\tau_{a,b}\int_a^b(t-a)t|t|\,dt,\nonumber\\
 L&=\sum_{(a,b)}\tau_{a,b}\int_{[a,b)}(t-a)\,d\nu(t).
 \label{ms:fourier:afcore:eq:gap-definitions}
\end{align}
The sums run over the same bounded complementary intervals of the closed
support as in $J$. The half-open interval is essential: an atom at the
upper endpoint is not an interior crossing of the gap.
If the support lies in $[-R,R]$, then
\[
 |J_1|\le RJ/2,\qquad |J_2|\le R^2J/2,
 \qquad 0\le L\le2R.
\]
For the first two bounds use $|t|\le R$ and
$\int_a^b(t-a)\,dt=(b-a)^2/2$. For the last, use
$\tau_{a,b}\le\E|X|\le1$, $t-a\le2R$, and disjointness of the
half-open gaps. Together with $J\le2R$, these bounds prove absolute
convergence of every sum.

\subsubsection{Integrating the same slack in two ways}

Put $\mathcal A=\int\tau(y)\{Q_H(y)-h_P(y)\}\,dy\ge0$.
On the support the slack vanishes; outside the convex hull $\tau=0$.
On a gap, contact at its endpoints gives
\[
 \nu([a,b))=h_P(a)-h_P(b)=Q_H(a)-Q_H(b)
                         =\int_a^b -Q_H'(t)\,dt.
\]
Bounded Fubini consequently gives
\[
 \int_a^b(Q_H-h_P)
 =\int_a^b(t-a)(-Q_H'(t))\,dt
       -\int_{[a,b)}(t-a)\,d\nu(t).
\]
In particular a mass at the upper endpoint is not counted as interior.
Substituting \eqref{ms:fourier:afcore:eq:contact-polynomial} proves
\begin{equation}\label{ms:fourier:afcore:eq:first-slack}
 \mathcal A=\frac{\varphi J}{2\sqrt n}+\frac{G}{n^{3/2}}-L,
 \quad
 G=\frac32C\beta J+(H\varphi+3C\rho)J_1-3CJ_2.
\end{equation}

The density $\tau$ is a probability density. If $X^*$ has that density,
signwise Fubini proves, for every $C^1$ test function on the compact range,
\[
 \E g'(X^*)=\E[Xg(X)],\qquad
 \E X^*=\gamma/2,\quad \E(X^*)^2=M_4/3,
 \quad\E|X^*|^3=M_5/4.
\]
Indeed, after subtracting $g(0)$, integration from zero to $X$ gives
$\E[Xg(X)]=\int\tau(t)g'(t)\,dt$ by signwise Fubini. This proves the
identity and the three moment formulas, using the relevant polynomial
or absolute-power test function.
Take $X^*$ independent of $S_{n-1}$. For a continuously differentiable
$g$, conditional application of that identity gives
\[
 \E g'((S_{n-1}+X^*)/\sqrt n)
 =\sqrt n\,\E[X_n g(S_n/\sqrt n)]
 =\E[Zg(Z)],\qquad Z=S_n/\sqrt n.
\]
The last equality is exchangeability. Hence
$Z^*=(S_{n-1}+X^*)/\sqrt n$ has the zero-bias density
$\tau_Z(t)=\E[Z\ind_{\{Z\ge t\}}]$, where $Z=S_n/\sqrt n$.
Indeed the same signwise Fubini formula for $Z$ identifies its integral
against every continuous test function, by taking that function as $g'$.
Its CDF at $z$ is
\[
 T=\E[(Z^2-zZ)\ind_{\{Z\le z\}}].
\]
For example, this formula follows by signwise integration of $\tau_Z$;
it is unaffected by endpoint atoms. Since $\E h_P(X^*)=T$, integration
of the contact polynomial against $\tau$ gives
\begin{equation}\label{ms:fourier:afcore:eq:second-slack}
 \mathcal A=F-\frac{\gamma\varphi}{2\sqrt n}
 +\frac{H\varphi(3-M_4)/6+CK_q}{n^{3/2}}-T,
 \quad K_q=\frac{M_5}{4}-\frac{\rho M_4}{2}
                  -\frac{3\beta\gamma}{2}+\frac\rho2.
\end{equation}

There is one further contact equality. Multiplication by $X$ and
exchangeability give $n\E[Xh_P(X)]=\E[S_n\ind_{\{S_n\le H\}}]$.
Thus, if
\[
 M=\E[Z\ind_{\{Z\le z\}}],\qquad
 K_1=\E[X|X|^3]-3\beta-\frac32\rho\gamma,
\]
then
\begin{equation}\label{ms:fourier:afcore:eq:truncated-first}
 M+\varphi=-\frac{\gamma z\varphi}{2\sqrt n}+\frac{CK_1}{n}.
\end{equation}
All moments here are legitimate because support was established first.

\subsubsection{Compensated inversion over the entire half-line}

\begin{lemma}[Quadratic hinge inversion]\label{ms:fourier:afcore:lem:inversion}
Let $\mu$ be a finite signed measure on $\mathbb R$ with zero mass,
zero first moment, and finite absolute second moment. With
$\widehat\mu(t)=\int e^{itx}\,d\mu(x)$, for every real $H$,
\begin{equation}\label{ms:fourier:afcore:eq:inversion}
 \int(x-H)|x-H|\,d\mu(x)
 =-\frac4\pi\int_0^\infty
        \frac{\Im(e^{-itH}\widehat\mu(t))}{t^3}\,dt.
\end{equation}
The right side is absolutely convergent.
\end{lemma}
\begin{proof}
Scaling and integration by parts give
\[
 \int_0^\infty\frac{ty-\sin(ty)}{t^3}\,dt
 =\frac\pi4 y|y|.
\]
For $y\ge0$ the integrand is nonnegative, and for $y<0$ it has the
opposite sign. Its absolute integral is therefore $\pi y^2/4$.
To check the constant, one integration by parts reduces it to half of
$\int_0^\infty(1-\cos t)t^{-2}\,dt$. Since
$t^{-2}=\int_0^\infty a e^{-at}\,da$, nonnegative Tonelli and elementary
Laplace integration give
\[
 \int_0^\infty\frac{1-\cos t}{t^2}\,dt
 =\int_0^\infty a\left\{\frac1a-\frac a{a^2+1}\right\}\,da
 =\int_0^\infty\frac{da}{a^2+1}=\frac\pi2.
\]
Absolute Fubini applies to the compensated integrand against $|\mu|$.
Its iterated absolute integral is
$\pi\int(x-H)^2\,d|\mu|(x)/4<\infty$.
The linear term integrates to zero, since both the mass and first moment
vanish. This proves the formula and its absolute convergence. No CDF
inversion or averaging of a jump is used.
\end{proof}

Let $Z_0$ be standard normal, write $f(t)=\E e^{itX}$, and set
\[
 q_z(x)=(x-z)|x-z|,\qquad
 \varepsilon=\E q_z(Z)-\E q_z(Z_0)-\frac{2\gamma\varphi}{3\sqrt n}.
\]
Apply Lemma~\ref{ms:fourier:afcore:lem:inversion} to the law of $S_n$ minus the
scale-$\sqrt n$ pushforward of the signed density
\[
 \phi(x)\left\{1+\frac{\gamma}{6\sqrt n}(x^3-3x)\right\}.
\]
That density has mass one and first moment zero. Gaussian integration
by parts gives its pushed-forward transform
$e^{-nt^2/2}(1-in\gamma t^3/6)$.
The hinge correction is $2\gamma\varphi/(3\sqrt n)$, because
$(x^3-3x)\phi(x)=-\phi'''(x)$ and
$-\int q_z\phi'''=4\phi(z)$, with the third derivative of $q_z$
understood distributionally as $4\delta_z$.
Since $q_H(\sqrt n\,x)=nq_z(x)$, the exact normalization is
\begin{equation}\label{ms:fourier:afcore:eq:epsilon}
 \varepsilon=-\frac4{\pi n}\int_0^\infty
        \frac{\Im(e^{-itH}V_n(t))}{t^3}\,dt,
 \quad V_n(t)=f(t)^n-e^{-nt^2/2}(1-in\gamma t^3/6).
\end{equation}
In particular the infinite high-frequency part has not been discarded.

\subsubsection{The identity and the error ledger}

\begin{proposition}[Conserved identity]\label{ms:fourier:afcore:prop:conserved}
Under the explicit contact hypotheses
\eqref{ms:fourier:afcore:eq:contact-input}--\eqref{ms:fourier:afcore:eq:contact-polynomial},
put $e=C-\CE$, $w=H^2/(2n)$, and
$U=J_1+(M_4-3)/6$. Then
\begin{align}
 ne\rho+n\CE\rho(1-e^{-w})
  +\frac{n\phi(0)e^{-w}}6D+n^{3/2}L
 ={}&C\{3\beta J/2+3\rho J_1-3J_2-K_q\}\nonumber\\
 &+H(CK_1+\varphi U)
   -H^2(C\rho+\gamma\varphi/2)+E_n,
 \label{ms:fourier:afcore:eq:conserved}\\
 E_n={}&-n^{3/2}\varepsilon/2.\nonumber
\end{align}
If $C\ge\CE$, every term on the left is nonnegative.
For any split $b>0$, valid complete bounds
\[
 I_{\rm low}\ge\int_0^b
       \frac{|\Im(e^{-itH}V_n(t))|}{t^3}\,dt,\qquad
 I_{\rm high}\ge\int_b^\infty\frac{|V_n(t)|}{t^3}\,dt
\]
give
\begin{equation}\label{ms:fourier:afcore:eq:error-ledger}
 |E_n|\le\frac{2\sqrt n}{\pi}(I_{\rm low}+I_{\rm high}).
\end{equation}
\end{proposition}
\begin{proof}
The pointwise identity
$q_z(x)=(x-z)^2-2(x-z)^2\ind_{\{x\le z\}}$ gives
\[
 T-\Phi(z)=-\tfrac12\{\E q_z(Z)-\E q_z(Z_0)\}
             +z(M+\varphi)-z^2(F-\Phi(z)).
\]
Here $\E[Z_0\ind_{\{Z_0\le z\}}]=-\varphi$ and the normal counterpart
of $T$ is $\Phi(z)$, by one integration by parts.
Use \eqref{ms:fourier:afcore:eq:truncated-first}, the definition of $\varepsilon$,
and $F-\Phi(z)=C\rho/\sqrt n$ in \eqref{ms:fourier:afcore:eq:second-slack}.
If
\[
 P_0=H^2(C\rho+\gamma\varphi/2)
    -H\{CK_1+\varphi(M_4-3)/6\}+CK_q,
\]
the result is
\[
 \sqrt n\mathcal A=C\rho-\gamma\varphi/6+P_0/n+\sqrt n\varepsilon/2.
\]
Comparison with \eqref{ms:fourier:afcore:eq:first-slack} yields
\[
 C\rho+\sqrt n\,L=\frac{\varphi}{6}(\gamma+3J)
                     +\frac{G-P_0}{n}-\sqrt n\varepsilon/2.
\]
Finally substitute
$\gamma+3J=(3+\sqrt{10})\rho-D$,
$\CE=(3+\sqrt{10})\phi(0)/6$, and
$\varphi=\phi(0)e^{-w}$, and multiply by $n$.
This proves \eqref{ms:fourier:afcore:eq:conserved}. Nonnegativity follows from
Lemma~\ref{ms:physical-defect}, $L\ge0$, and $e\ge0$.
Equation~\eqref{ms:fourier:afcore:eq:epsilon} gives the ledger directly.
\end{proof}

The identity holds at each actual index $n\ge2$. When $C\ge\CE$,
its four nonnegative terms have a direct interpretation: score excess, Gaussian phase cost,
physical defect, and leakage into complementary support gaps. The next
step turns a complete bound on the right side into a quantitative budget.

\subsubsection{Removing the Gaussian weight and retaining its phase cost}

Define the exact raw remainder, without subtracting a score charge, by
\begin{equation}\label{ms:fourier:rescore:re-raw}
\begin{aligned}
 \mathcal R={}&C\{3\beta J/2+3\rho J_1-3J_2-K_q\}
 +H(CK_1+\varphi U)\\
 &-H^2(C\rho+\gamma\varphi/2)+E_n,\qquad
 U=J_1+(M_4-3)/6.
\end{aligned}
\end{equation}

The exact conserved identity can be regrouped as
\begin{equation}\label{ms:conserved:regrouped}
 \mathcal R=ne\rho+\frac{n\phi(0)}6D+n^{3/2}L
 +n(1-e^{-w})\left(\CE\rho-\frac{\phi(0)}6D\right).
\end{equation}
The last coefficient must be positive before the Gaussian weight can
be removed. A first complete remainder bound provides that sign.

\begin{lemma}[Removing the Gaussian weight]\label{ms:fourier:rescore:re-weight}
Suppose $C\ge\CE$, $n\ge N\ge2$, and an established bound gives
$\mathcal R\le r_0<2N/5$. Then
\begin{equation}\label{ms:fourier:rescore:re-unweighted}
 ne\rho+\frac{n\phi(0)}6D+n^{3/2}L\le\mathcal R.
\end{equation}
Consequently any later bound $\mathcal R\le r$ gives
\[
 e^{-w}\ge1-\frac{5r}{2n},\qquad
 nD\le\frac{15000r}{997}.
\]
\end{lemma}
\begin{proof}
The identity before substitution of $D$ gives
\[
 \frac{\varphi}{6}(\gamma+3J)
 =C\rho+\sqrt n\,L-\frac{\mathcal R}{n}
 >C\rho-\frac25>0,
\]
using $\rho\ge1$ and $C\ge\CE>2/5$. Thus
$\CE\rho-\phi(0)D/6=\phi(0)(\gamma+3J)/6>0$.
Dropping the last nonnegative term of
\eqref{ms:conserved:regrouped} proves
\eqref{ms:fourier:rescore:re-unweighted}.
In the original conserved identity all four terms are nonnegative,
so $n\CE\rho(1-e^{-w})\le r$. Since $\CE\rho>2/5$, this proves
the density floor. Finally $\phi(0)>997/2500$ gives the bound on $nD$.
If a purported upper bound has $r<0$, the same inequalities already
exclude the parameter cell.
\end{proof}

The constant comparisons used here follow from
Lemma~\ref{ms:gaussian-constants}.

\begin{corollary}[Conserved phase resource]\label{ms:conserved:phase-resource}
Under the preceding hypotheses, suppose established bounds give
\[
 \mathcal R\le r<2N/5,\qquad nD\le Nd_D,\qquad
 a_{\rm ph}(\rho):=0.4095\rho-d_D/15>0.
\]
Suppose also that on the current signed phase cell the raw remainder
has a pointwise quadratic envelope
\[
 \mathcal R(H)\le A(C,\rho)+F_{\operatorname{sgn}H}(C,\rho)|H|
                  -Q_{\rm raw}(C,\rho)H^2.
\]
Set
\begin{equation}\label{ms:fourier:rescore:re-extra}
 q_{\rm extra}(\rho)=\frac{1-5r/(2N)}2
                   \left(0.4095\rho-\frac{d_D}{15}\right)>0
\end{equation}
and
\[
 \mathcal B(H)=A(C,\rho)+F_{\operatorname{sgn}H}(C,\rho)|H|
                -\{Q_{\rm raw}(C,\rho)+q_{\rm extra}(\rho)\}H^2.
\]
Then, with $\kappa=997/15000$,
\begin{equation}\label{ms:fourier:rescore:re-phaseprice}
 n(C-\CE)\rho+\kappa nD+n^{3/2}L\le\mathcal B(H).
\end{equation}
The same assertion holds if the quadratic envelope is a maximum of
several valid corner quadratics, after adding the phase price to each
corresponding curvature.
\end{corollary}
\begin{proof}
Because $n\ge N$, the weighted defect bound implies $D\le d_D$.
The sign assumptions, $\CE>0.4095$ and $\phi(0)<2/5$ yield
\[
 \CE\rho-\phi(0)D/6\ge0.4095\rho-d_D/15>0.
\]
Also $e^{-w}\ge1-5r/(2n)\ge1-5r/(2N)>0$.
For $w\ge0$, the inequality $e^w\ge1+w$ gives
$1-e^{-w}\ge we^{-w}$. The last term in
\eqref{ms:conserved:regrouped} is therefore at least
$q_{\rm extra}(\rho)H^2$. Subtract this term from the raw upper bound
and use $\phi(0)/6>\kappa$ to obtain
\eqref{ms:fourier:rescore:re-phaseprice}.
\end{proof}

This corollary specifies how the Fourier remainder enters the resource
calculation. Complete low- and high-frequency bounds in
Appendix~\ref{ms:fourier:rescore:sec:setup} control $E_n$; moment and gap
bounds control the other terms of $\mathcal R$. The coefficient atlas in
Appendix~\ref{ms:fourier:rescore:re-atlas} organizes these enclosures into
the signed quadratics above. Once a phase price has been included in
$Q_{\rm raw}+q_{\rm extra}$, it is not added again. The charge
$n(C-\CE)\rho$ remains on the left throughout, ready for the convex
score argument that follows.
